% Source: source/普通高考/2010/2010辽宁理.pdf, page 1, question 4.
% Complete node -> directed-edge list:
%   开始 -> 输入 n,m -> k=1,p=1 -> p=p(n-m+k) -> k<m;
%   是 -> k=k+1 -> left return to the shared stem above the update;
%   否 -> 输出 p -> 结束.
% The return edge is independent; the shared stem has one downward arrow into
% the update node and no duplicate arrow at the loop junction.

\begin{tikzpicture}[
  >=Stealth,
  x=1cm,
  y=0.82cm,
  line width=0.45pt,
  line cap=round,
  line join=round,
  font=\normalsize,
  flowtext/.style={anchor=center, align=center,
    execute at begin node={\setlength{\parindent}{0pt}\noindent}},
  every node/.style={flowtext},
  flowbox/.style={flowtext, draw, minimum height=0.68cm,
    inner xsep=4pt, inner ysep=2.5pt},
  branchlabel/.style={flowtext, inner sep=0pt},
  terminal/.style={flowbox, rounded corners=0.18cm, minimum width=1.35cm},
  process/.style={flowbox},
  decision/.style={flowbox, diamond, aspect=3.0, minimum width=3.45cm,
    minimum height=0.95cm},
  io/.style={flowbox, trapezium, trapezium left angle=72,
    trapezium right angle=108, minimum width=2.10cm, minimum height=0.76cm},
]
  \path[use as bounding box] (-2.10,0.05) rectangle (5.10,9.65);

  \node[terminal] (start) at (0,9.30) {开始};
  \node[io, minimum width=2.45cm] (input) at (0,8.00) {输入 $n,m$};
  \node[process, minimum width=2.60cm] (init) at (0,6.80)
    {$k=1,p=1$};
  \node[process, minimum width=3.55cm] (updateP) at (0,5.27)
    {$p=p(n-m+k)$};
  \node[decision] (test) at (0,3.85) {$k<m$};
  \node[process, minimum width=2.25cm] (inc) at (3.35,5.35)
    {$k=k+1$};
  \node[io] (output) at (0,1.95) {输出 $p$};
  \node[terminal] (finish) at (0,0.55) {结束};

  \draw[->] (start.south) -- (input.north);
  \draw[->] (input.south) -- (init.north);
  \coordinate (merge) at (0,6.00);
  \coordinate (branchright) at (3.35,3.85);
  \draw[->] (init.south) -- (updateP.north);
  \draw[->] (updateP.south) -- (test.north);

  % Yes branch: right, then upward into the increment node.
  \draw (test.east) --
    node[branchlabel, pos=.65, above=1pt] {是} (branchright);
  \draw[->] (branchright) -- (inc.south);
  \draw[->] (inc.north) -- (3.35,6.00) -- (merge);

  % No branch exits directly to output and termination.
  \draw[->] (test.south) -- node[branchlabel, right=3pt] {否} (output.north);
  \draw[->] (output.south) -- (finish.north);
\end{tikzpicture}
